% Source: https://artofproblemsolving.com/community/c3h2601591_probability_problem
% Stable identity: aops-topic-2601591-q01
The $2^z$ outcomes are equally likely, so any test has success probability $\dfrac{k}{2^z}$, where $k$ is the number of successful outcomes.

\textbf{(a)} If $\dfrac{k}{2^z}=\dfrac13$, then $3k=2^z$, which is impossible because $3$ does not divide a power of $2$.

\textbf{(b)} If $\dfrac xy=\dfrac{k}{2^z}$, then $x\cdot2^z=ky$. Since $\gcd(x,y)=1$, $y$ divides $2^z$, so $y$ is a power of $2$. Conversely, if $y=2^m$, toss $z=m$ coins and declare success for any $x$ of the $2^m$ outcomes. So the achievable probabilities are exactly those with $\boxed{y\text{ a power of }2}$. In particular, not every fraction $\dfrac xy$ is possible. $\blacksquare$
